% !TEX program = pdflatex
% Компиляция: pdflatex report_oil_gestures.tex (дважды, для оглавления и ссылок)
\documentclass[11pt,a4paper]{article}

% --- Кодировки и язык ---
\usepackage[T2A]{fontenc}
\usepackage[utf8]{inputenc}
\usepackage[russian]{babel}

% --- Геометрия и оформление ---
\usepackage[a4paper,margin=2.3cm,headheight=14pt]{geometry}
\usepackage{xcolor}
\usepackage{booktabs}
\usepackage{longtable}
\usepackage{array}
\usepackage{enumitem}
\usepackage{titlesec}
\usepackage{fancyhdr}
\usepackage{tikz}
\usetikzlibrary{shapes.geometric, arrows.meta, positioning, fit, backgrounds}
\usepackage{listings}
\usepackage{amsmath}
\usepackage{amssymb}
\usepackage[hidelinks]{hyperref}

% --- Цветовая палитра ---
\definecolor{accent}{RGB}{17,84,140}
\definecolor{accentlight}{RGB}{225,236,247}
\definecolor{softgray}{RGB}{96,96,96}
\definecolor{codebg}{RGB}{246,248,250}
\definecolor{rulec}{RGB}{200,210,220}

% --- Гиперссылки ---
\hypersetup{
  colorlinks=true,
  linkcolor=accent,
  urlcolor=accent,
  citecolor=accent,
  pdftitle={Oil Gestures — недельный отчёт о проделанной работе},
  pdfauthor={Mukhtashev Azamat Manuchehr}
}

% --- Колонтитулы ---
\pagestyle{fancy}
\fancyhf{}
\renewcommand{\headrulewidth}{0.4pt}
\renewcommand{\footrulewidth}{0pt}
\lhead{\small\textcolor{softgray}{Oil Gestures}}
\rhead{\small\textcolor{softgray}{Отчёт о проделанной работе}}
\cfoot{\thepage}

% --- Стиль заголовков ---
\titleformat{\section}
  {\color{accent}\Large\bfseries}{\thesection}{0.6em}{}
\titleformat{\subsection}
  {\color{accent!85!black}\large\bfseries}{\thesubsection}{0.6em}{}
\titlespacing*{\section}{0pt}{16pt}{8pt}

% --- Оформление листингов кода ---
\lstdefinestyle{mono}{
  backgroundcolor=\color{codebg},
  basicstyle=\ttfamily\footnotesize,
  breaklines=true,
  frame=single,
  rulecolor=\color{rulec},
  framesep=6pt,
  xleftmargin=6pt,
  xrightmargin=2pt,
  showstringspaces=false,
  keywordstyle=\color{accent}\bfseries,
  commentstyle=\color{softgray}\itshape,
  language=Python,
  literate={->}{{$\rightarrow$}}1
}
\lstset{style=mono}

% --- Вспомогательная команда для блока «итог» ---
\newcommand{\highlight}[1]{\colorbox{accentlight}{\parbox{\dimexpr\linewidth-2\fboxsep}{#1}}}

% =====================================================================
\begin{document}

% --------------------------- Титул -----------------------------------
\begin{titlepage}
\centering
\vspace*{1.5cm}
{\color{accent}\rule{\linewidth}{1.5pt}}\\[0.5cm]
{\Huge\bfseries Oil Gestures\par}
\vspace{0.4cm}
{\Large Система распознавания жестов руки для управления\\ демонстрационным симулятором нефтедобычи\par}
\vspace{0.3cm}
{\large\color{softgray} Отчёт о проделанной работе за неделю\par}
{\color{accent}\rule{\linewidth}{1.5pt}}\\[1.2cm]

\begin{tabular}{rl}
\textbf{Отчётный период:} & 16--22 июня 2026 г. \\[4pt]
\textbf{Исполнитель:} & Mukhtashev Azamat Manuchehr \\[4pt]
\textbf{Дата отчёта:} & 22 июня 2026 г. \\
\end{tabular}

\vspace{1.4cm}

\highlight{\textbf{Краткое резюме.} За отчётную неделю реализован сквозной конвейер распознавания
жестов на базе MediaPipe, добавлена изолированная подсистема управления курсором с
кроссплатформенными бэкендами мыши (macOS, Windows, Linux/X11, Linux/Wayland), оптимизирована
производительность на Linux и спроектирован версионируемый контракт интеграции ML с UI- и
3D-приложениями. Реализовано распознавание 4 статических и 4 курсорных жестов с производным
перетаскиванием; тестовый набор — \textbf{71 тест} (68 проходят, 3 пропущены по платформенным
условиям).}

\vfill
\end{titlepage}

\tableofcontents
\thispagestyle{fancy}
\newpage

% =====================================================================
\section{Цель и общие сведения о проекте}

\textbf{Oil Gestures} — система компьютерного зрения, которая по изображению с веб-камеры
распознаёт жесты руки и преобразует их в команды и события для демонстрационного симулятора
нефтедобычи. Ядром продукта является конвейер распознавания:

\begin{center}
\ttfamily Камера $\rightarrow$ ландмарки MediaPipe $\rightarrow$ независимые распознаватели $\rightarrow$ команды / события
\end{center}

MediaPipe выступает единственным источником координат ключевых точек руки (landmarks).
Управление курсором реализовано как \emph{вторичная} подсистема с собственным распознавателем:
она использует общий поток ландмарок, но не потребляет результаты статических или динамических
жестов как действия мыши. Такое разделение позволяет развивать распознавание жестов и
управление курсором независимо.

\subsection{Архитектура времени исполнения}

Система построена как слоистая архитектура: каждый слой решает одну задачу и обменивается
данными через простые типизированные структуры с явно описанными границами.

\begin{center}
\begin{tikzpicture}[
  node distance=8mm and 10mm,
  box/.style={rectangle, rounded corners, draw=accent, fill=accentlight,
              align=center, minimum height=9mm, inner sep=4pt, font=\small},
  sub/.style={rectangle, rounded corners, draw=rulec, fill=codebg,
              align=center, minimum height=8mm, inner sep=3pt, font=\footnotesize},
  arr/.style={-{Latex[length=2mm]}, draw=accent, thick}
]
\node[box] (cam) {Камера\\\scriptsize OpenCV / V4L2};
\node[box, right=of cam] (proc) {Обработка кадра\\\scriptsize mirror / resize / RGB};
\node[box, right=of proc] (mp) {MediaPipe\\\scriptsize landmarker};

\node[sub, below=14mm of cam] (static) {Статические\\жесты};
\node[sub, right=of static] (dyn) {Динамические\\(модель)};
\node[sub, right=of dyn] (cur) {Курсорный\\распознаватель};

\node[box, below=10mm of dyn] (pipe) {Курсорный конвейер\\\scriptsize сглаживание + бэкенд мыши};
\node[box, right=22mm of pipe] (pub) {Публикатор ML\\\scriptsize поток событий по сети};

\draw[arr] (cam) -- (proc);
\draw[arr] (proc) -- (mp);
\draw[arr] (mp.south) .. controls +(0,-6mm) and +(0,6mm) .. (static.north);
\draw[arr] (mp.south) -- (dyn.north);
\draw[arr] (mp.south) .. controls +(0,-6mm) and +(0,6mm) .. (cur.north);
\draw[arr] (cur) -- (pipe);
\draw[arr] (static.south) .. controls +(0,-6mm) and +(0,6mm) .. (pub.north west);
\draw[arr] (pipe) -- (pub);
\end{tikzpicture}
\end{center}

\noindent Композиция модулей и цикл OpenCV сосредоточены в \texttt{app/main.py}; ни один
жестовый алгоритм не находится в слое приложения — он лишь связывает компоненты.

\subsection{Принципы проектирования}

В основу архитектуры положены несколько идей, выдержанных по всему проекту:

\begin{itemize}[leftmargin=1.4em]
  \item \textbf{Чёткие границы между слоями.} Слои обмениваются простыми пакетами данных
        (кадр, ландмарки руки, результат жеста, позиция указателя и т.\,п.) с явно описанными
        «производителем» и «потребителями», поэтому подсистемы можно развивать независимо.
  \item \textbf{Независимые распознаватели.} Статические, динамические и курсорные жесты
        распознаются отдельными модулями; вместе с каждым результатом передаётся его источник,
        что исключает смешивание подсистем.
  \item \textbf{Изоляция работы с операционной системой.} Управление мышью спрятано за единым
        интерфейсом, а платформенные реализации подключаются только под нужную ОС. Благодаря
        этому остальной код не зависит от конкретной платформы, а тесты идут на любой системе.
  \item \textbf{Конфигурация как данные.} Поведение задаётся YAML-настройками и флагами командной
        строки, а не правками кода.
\end{itemize}

\subsection{Структура репозитория}

\begin{center}
\begin{tabular}{>{\ttfamily}l l}
\toprule
\normalfont\textbf{Каталог} & \textbf{Назначение} \\
\midrule
app/                 & Типизированная конфигурация и цикл исполнения \\
oil\_gestures/vision/  & Камера, MediaPipe, отрисовка, работа с ландмарками \\
oil\_gestures/gestures/ & Независимые распознаватели: статика, динамика, курсор \\
oil\_gestures/cursor/   & Указатель, сглаживание, фасад и бэкенды мыши \\
oil\_gestures/integration/ & Контракты ML, сервер событий, публикатор, клиент \\
oil\_gestures/commands/ & Сопоставление жестов и команд симулятора \\
oil\_gestures/core/     & Перечисления, типы, константы, логирование \\
configs/             & YAML-конфигурация по умолчанию \\
contracts/           & Версионируемая JSON-схема событий ML \\
scripts/             & Тонкие исполняемые скрипты \\
tests/               & Модульные тесты чистой логики \\
docs/                & Проектная документация \\
\bottomrule
\end{tabular}
\end{center}

% =====================================================================
\section{Выполненные работы за неделю}

Ниже описаны реализованные подсистемы: что делает каждая часть и как она работает на
концептуальном уровне.

\subsection{Конвейер компьютерного зрения и захват камеры}

\begin{itemize}[leftmargin=1.4em]
  \item \textbf{Единый анализ кадра.} За один проход MediaPipe система получает сразу и точки
        руки (для курсора и динамики), и встроенную категорию жеста (для статики). Это
        вдвое экономит вычисления по сравнению с запуском двух отдельных моделей. Файл модели при
        отсутствии скачивается автоматически.
  \item \textbf{Два разрешения кадра.} Камера снимает в высоком разрешении; из него готовятся
        отдельно уменьшенный кадр для распознавания и кадр для показа на экране. Так точность
        детекции сохраняется, а отрисовка предпросмотра не нагружает систему.
  \item \textbf{Захват в отдельном потоке.} Чтение с камеры идёт параллельно с распознаванием и
        всегда отдаёт самый свежий кадр: если обработка не успевает, устаревшие кадры
        пропускаются и задержка не накапливается.
  \item \textbf{Подстройка под камеру и ОС.} Автоматически выбираются подходящий способ захвата и
        формат потока для текущей системы; фактические параметры камеры (разрешение, частота,
        формат) логируются, а при низкой частоте кадров выводится предупреждение с рекомендацией.
  \item \textbf{Метрики производительности.} Частота кадров и время распознавания измеряются и
        сглаживаются, что даёт стабильные показатели для диагностики.
\end{itemize}

\subsection{Распознавание статических жестов}

Статические жесты распознаются на основе встроенного классификатора MediaPipe: используется
категория, уже посчитанная при едином анализе кадра, поэтому отдельная нейросеть или ручные
геометрические правила не нужны. Слишком неуверенные результаты отбрасываются, а из всех
категорий MediaPipe берутся только нужные проекту — остальные намеренно игнорируются, что
снижает число ложных срабатываний. Соответствие категорий проектным жестам:

\begin{center}
\begin{tabular}{>{\ttfamily}l c >{\ttfamily}l}
\toprule
\normalfont\textbf{Категория MediaPipe} & $\rightarrow$ & \normalfont\textbf{Жест проекта} \\
\midrule
Closed\_Fist & $\rightarrow$ & FIST \\
Open\_Palm   & $\rightarrow$ & OPEN\_PALM \\
Thumb\_Up    & $\rightarrow$ & THUMB\_UP \\
Victory      & $\rightarrow$ & VICTORY \\
\bottomrule
\end{tabular}
\end{center}

\noindent Результаты фильтруются по порогу уверенности; нераспознанные или низкоуверенные
категории отбрасываются. Распознаватель не содержит правил курсора, что сохраняет границы
подсистем.

\subsection{Подсистема управления курсором}

Изолированная вторичная подсистема, не использующая результаты статических и динамических жестов.
Её работа складывается из нескольких шагов:

\begin{itemize}[leftmargin=1.4em]
  \item \textbf{Распознавание жестов курсора.} Курсорные команды определяются по взаимному
        положению пальцев (например, сближение большого и указательного пальцев — «щипок»).
        Расстояние нормируется по размеру кисти, поэтому жест распознаётся одинаково независимо от
        того, ближе или дальше рука от камеры. Чтобы жест не «дребезжал» на границе срабатывания,
        используются разные пороги на нажатие и на отпускание.
  \item \textbf{Указатель и перенос на экран.} За положение курсора берётся одна точка руки и
        переносится в координаты экрана. По краям кадра оставлены поля (активная зона), чтобы до
        любого угла экрана можно было дотянуться без выхода руки из поля зрения; поддерживаются
        несколько мониторов.
  \item \textbf{Сглаживание движения.} Координаты курсора проходят через адаптивный сглаживающий
        фильтр: при медленных движениях он сильно подавляет дрожание руки, а при быстрых почти не
        вносит задержки. Дополнительно отсекаются микродвижения и ограничивается максимальная
        скорость скачка.
  \item \textbf{Логика управления.} Управление включается только после устойчивого захвата руки,
        что исключает случайные скачки при кратковременной потере трекинга. Поддержаны клик,
        захват/отпускание и перетаскивание; при потере руки удержанная кнопка безопасно
        отпускается, а повторные срабатывания одного действия ограничены небольшой паузой.
  \item \textbf{Включение режима жестом.} Курсорный режим включается и выключается удержанием
        статического жеста. Реализована защита от ложных срабатываний: жест нужно удержать
        некоторое время, затем отпустить, а между переключениями выдерживается минимальный
        интервал; на экран выводится индикатор прогресса удержания.
  \item \textbf{Безопасный режим по умолчанию.} Реальная мышь не двигается, пока это не разрешено
        явно (флаг \texttt{--real-mouse}); на macOS дополнительно проверяется системное разрешение
        с понятной подсказкой.
\end{itemize}

\subsection{Кроссплатформенные бэкенды мыши}

Управление мышью спрятано за единым интерфейсом, а конкретная реализация выбирается под текущую
операционную систему и подключается только при необходимости. Это позволяет поддерживать разные
платформы, не усложняя остальной код.

\begin{center}
\begin{tabular}{l l l}
\toprule
\textbf{Платформа} & \textbf{Способ управления} & \textbf{Особенности} \\
\midrule
macOS          & Системный API курсора   & Проверка/запрос разрешения \\
Windows        & Системный API курсора   & Перемещение и клики \\
Linux / X11    & Стандартный механизм X11 & Сессии X11/Xorg \\
Linux / Wayland & Виртуальное устройство ядра & Реальный курсор под Wayland \\
Любая (тест)   & Безопасная заглушка     & Без обращения к ОС \\
\bottomrule
\end{tabular}
\end{center}

\noindent Для каждой платформы используется её штатный способ управления указателем, а для
теста и безопасного режима — заглушка, которая ведёт позицию курсора, но не трогает реальную
систему.

\highlight{\textbf{Ключевое достижение недели:} реализована поддержка реального курсора под
\emph{Wayland} на Linux. Привычный для X11 способ синтезировать события мыши под Wayland не
двигает настоящий курсор, поэтому был добавлен отдельный бэкенд, который создаёт виртуальное
указательное устройство на уровне ядра — так же, как его «видела» бы система от физического
устройства. Благодаря этому курсор корректно двигается и в современных Wayland-сессиях; при
отсутствии нужных прав доступа пользователь получает понятное сообщение с инструкцией.}

\subsection{Оптимизация производительности на Linux}

\begin{itemize}[leftmargin=1.4em]
  \item Для камеры по умолчанию запрашивается сжатый формат потока, что снимает ограничение
        пропускной способности и повышает частоту кадров.
  \item Уменьшение кадра предпросмотра перед показом вынесено в настройки; это устранило
        основную статью затрат производительности на Linux, не затрагивая качество распознавания.
  \item В журнал выводятся фактические параметры камеры — для быстрой диагностики
        производительности на разных машинах.
\end{itemize}

\subsection{Контракт автономной интеграции ML}

Спроектирован контракт интеграции, по которому модуль распознавания работает как самостоятельный
\emph{поставщик} событий, а приложения интерфейса и 3D-сцены — как независимые \emph{потребители}.
Стороны не зависят от внутреннего кода друг друга; их единственная общая договорённость — формат
передаваемых данных.

\begin{itemize}[leftmargin=1.4em]
  \item \textbf{Передача данных.} События передаются по локальному сетевому соединению в виде
        простых текстовых JSON-сообщений. К поставщику может одновременно подключаться несколько
        приложений, при этом модуль распознавания никогда не ждёт их: если приложение не успевает
        читать, оно просто получит самые свежие данные, а не накопившуюся очередь.
  \item \textbf{Состояние на каждом кадре.} В каждом сообщении передаётся полная картина: параметры
        кадра, частота кадров, обнаружена ли рука и где она находится, результаты всех
        распознавателей (статика, динамика, курсор), состояние курсорного режима и паузы. Каждое
        сообщение пронумеровано, что позволяет приложениям отслеживать порядок и не путать данные.
  \item \textbf{Опциональная картинка с камеры.} По запросу поставщик может дополнительно
        отправлять кадры с камеры с регулируемой частотой и качеством, чтобы приложения могли
        показывать видео и рисовать поверх него свои подсказки. Это не влияет на скорость
        распознавания.
  \item \textbf{Расширяемость без поломок.} Формат рассчитан на развитие: новые жесты можно
        добавлять, не нарушая работу уже подключённых приложений.
  \item \textbf{Формальное описание.} Структура сообщений зафиксирована в виде схемы и проверяется
        автотестами; приложен готовый пример приложения-потребителя.
\end{itemize}

\begin{lstlisting}[caption={Запуск автономного поставщика событий},captionpos=b]
# Поставщик: данные распознавания + кадры камеры, без локального окна
python scripts/run_demo.py --event-server --publish-camera --headless
\end{lstlisting}

\subsection{Конфигурация, командный интерфейс и задел на динамику}

\begin{itemize}[leftmargin=1.4em]
  \item \textbf{Настройка через файлы и флаги.} Поведение системы (камера, режимы распознавания,
        курсор, действия) задаётся YAML-настройками, которые при необходимости переопределяются
        флагами командной строки.
  \item \textbf{Удобные режимы запуска и диагностики.} Добавлены флаги для включения/выключения
        курсора, разрешения реального управления мышью, headless-режима без окна и публикации
        событий, а также диагностические режимы, проверяющие управление мышью и параметры экрана
        без запуска камеры.
  \item \textbf{Задел под динамические жесты.} Подготовлен интерфейс для будущей обученной модели
        динамических жестов: он принимает результаты только от такой модели и накапливает короткую
        историю движений руки. Это заранее фиксирует границу подключения модели, не требуя в
        дальнейшем переделки остальной системы.
\end{itemize}

\subsection{Сводка ключевых механизмов}

\begin{center}
\begin{longtable}{l p{8.0cm}}
\toprule
\textbf{Механизм} & \textbf{Что даёт} \\
\midrule
\endhead
Единый анализ кадра & Точки руки и жест за один проход — вдвое меньше вычислений \\
Захват в отдельном потоке & Всегда свежий кадр, задержка не накапливается \\
Два разрешения кадра & Точная детекция при недорогом предпросмотре \\
Нормировка жеста по размеру кисти & Распознавание не зависит от удаления руки \\
Два порога на жест (нажатие/отпускание) & Устойчивость без «дребезга» на границе \\
Адаптивное сглаживание курсора & Нет дрожания на медленных и задержки на быстрых движениях \\
Активная зона и поддержка мониторов & Доступ ко всему экрану без выхода руки из кадра \\
Логика захвата и пауз & Защита от случайных скачков и «залипших» кнопок \\
Защищённое включение режима жестом & Нет ложных переключений \\
Единый интерфейс к мыши & Поддержка macOS, Windows, Linux без усложнения кода \\
Виртуальное устройство для Wayland & Реальный курсор в современных Linux-сессиях \\
Неблокирующая раздача событий & Несколько приложений-потребителей не тормозят распознавание \\
Зафиксированный формат данных & Стабильная и расширяемая граница с UI и 3D \\
\bottomrule
\end{longtable}
\end{center}

% =====================================================================
\section{Тестирование и контроль качества}

Создан набор модульных тестов чистой логики (без зависимости от камеры и ОС). Текущее
состояние прогона:

\begin{center}
\highlight{\centering\textbf{71 тест собран $\;\bullet\;$ 68 пройдено $\;\bullet\;$ 3 пропущено}
(пропуски обусловлены платформенными условиями), время прогона $\approx$\,1,1\,с.}
\end{center}

\begin{center}
\begin{tabular}{>{\ttfamily}l c l}
\toprule
\normalfont\textbf{Файл тестов} & \textbf{Тестов} & \textbf{Покрытая область} \\
\midrule
test\_cursor\_gestures.py      & 18 & Распознавание курсорных жестов, гистерезис \\
test\_platform\_backends.py    & 11 & Выбор и контракт бэкендов мыши \\
test\_gesture\_toggle.py       & 10 & Переключение режима удержанием жеста \\
test\_integration\_contracts.py & 7 & Сериализация и схема событий ML \\
test\_static\_recognizer.py    & 5  & Отображение canned-жестов MediaPipe \\
test\_vision\_performance.py   & 4  & Формат камеры, потоковый захват, кадр инференса \\
test\_cooldown.py             & 3  & Перезарядка действий \\
test\_command\_mapper.py       & 2  & Сопоставление жестов и команд \\
test\_cursor\_boundaries.py    & 2  & Границы отображения в экран \\
test\_sequence\_buffer.py      & 2  & Буфер последовательностей (динамика) \\
test\_runtime.py              & 1  & Композиция рантайма \\
\bottomrule
\end{tabular}
\end{center}

% =====================================================================
\section{Технические показатели}

\begin{center}
\begin{tabular}{l r}
\toprule
\textbf{Показатель} & \textbf{Значение} \\
\midrule
Объём кода Python (без \texttt{.venv}) & $\approx$ 6\,400 строк \\
Реализованных модулей Python & 67 \\
Поддерживаемых статических жестов & 4 \\
Поддерживаемых курсорных жестов & 4 (+ перетаскивание) \\
Модульных тестов & 71 (68 проходят) \\
Платформенных бэкендов мыши & 5 \\
Документов в \texttt{docs/} & 9 \\
Целевой рантайм & CPython 3.14.6 \\
Версия контракта ML & v1 (JSON Schema Draft 2020-12) \\
\bottomrule
\end{tabular}
\end{center}

% =====================================================================
\section{Использованные инструменты и их версии}

Все зависимости зафиксированы по версиям для воспроизводимости сборки. Целевой
рантайм — стандартная сборка CPython 3.14.6 (свободно-поточная сборка \texttt{3.14t}
пока не поддерживается).

\subsection{Среда исполнения и основные библиотеки}

\begin{center}
\begin{tabular}{l >{\ttfamily}l c}
\toprule
\textbf{Назначение} & \normalfont\textbf{Инструмент} & \textbf{Версия} \\
\midrule
Язык и рантайм               & CPython                     & 3.14.6 \\
Распознавание руки (landmarks) & mediapipe                 & 0.10.35 \\
Обработка изображений и захват & opencv-contrib-python      & 4.13.0.92 \\
Численные вычисления          & numpy                       & 2.4.6 \\
Загрузка конфигурации (YAML)  & PyYAML                      & 6.0.3 \\
\bottomrule
\end{tabular}
\end{center}

\subsection{Платформенные бэкенды мыши}

\begin{center}
\begin{tabular}{l >{\ttfamily}l c}
\toprule
\textbf{Платформа} & \normalfont\textbf{Инструмент} & \textbf{Версия} \\
\midrule
macOS          & pyobjc-framework-Quartz & 11.1 \\
Linux / X11    & python-xlib             & 0.33 \\
Linux / Wayland & \normalfont ядерный \texttt{/dev/uinput} (стандартная библиотека) & --- \\
Windows        & \normalfont WinAPI \texttt{ctypes} (стандартная библиотека) & --- \\
\bottomrule
\end{tabular}
\end{center}

\subsection{Инструменты разработки и контроля качества}

\begin{center}
\begin{tabular}{l >{\ttfamily}l c}
\toprule
\textbf{Назначение} & \normalfont\textbf{Инструмент} & \textbf{Версия} \\
\midrule
Модульное тестирование        & pytest      & 8.4.2 \\
Валидация JSON-схемы контракта & jsonschema  & 4.25.1 \\
Изоляция окружения            & \normalfont venv (стандартная библиотека) & --- \\
Контроль версий               & \normalfont Git & --- \\
\bottomrule
\end{tabular}
\end{center}

% =====================================================================
\section{Результаты и текущий статус}

\begin{itemize}[leftmargin=1.4em]
  \item Сквозной конвейер «камера $\rightarrow$ MediaPipe $\rightarrow$ распознаватели» работает и
        запускается одной командой (\texttt{python scripts/run\_demo.py}).
  \item Статическое распознавание поддерживает 4 жеста; курсорное — 4 жеста с производным
        перетаскиванием.
  \item Реальное управление курсором подтверждено на macOS, Windows, Linux/X11 и
        Linux/Wayland; безопасный dry-run включён по умолчанию.
  \item Контракт интеграции ML формализован схемой и готов к независимому подключению UI- и
        3D-приложений.
  \item Качество зафиксировано модульными тестами (68 проходят) и проектной документацией.
\end{itemize}

% =====================================================================
\section{Дальнейшие шаги}

\begin{enumerate}[leftmargin=1.6em]
  \item Подключить обученную модель динамических жестов к подготовленному
        \emph{model-only} интерфейсу (pointing, OK, swipe, clench/release, поворот запястья).
  \item Реализовать автоматическую активацию режима курсора обученным динамическим жестом на
        уровне оркестрации рантайма.
  \item Расширить сопоставление жестов и команд симулятора (давление, клапаны, насосы,
        аварийная остановка) поверх существующего перечня \texttt{CommandName}.
  \item Подключить реальные UI- и 3D-приложения как независимых потребителей контракта v1 и
        провести нагрузочную проверку нескольких одновременных потребителей.
\end{enumerate}

\vspace{0.8cm}
\noindent\textcolor{rulec}{\rule{\linewidth}{0.6pt}}\\[4pt]
{\small\color{softgray} Целевой рантайм — CPython 3.14.6. Запуск демонстрации:
\texttt{python scripts/run\_demo.py}.}

\end{document}
